%% kapa: exam=1 round=31 year=2020 period=2 form=A subject=appraisal_law subject_ko=감정평가관계법규 seq=8 number=8
%% source: Q-Net 감정평가사 기출문제 (cst00309) artlSeq=5207219 file=extracted/1008/5207219/A형/A형/2020년도 제31회 감정평가사 1차 시험_2교시_A형.pdf
%% conversion: kordoc 4.24.51 -> tools/gen1.py (auto); review=proofread-sonnet; flags=none
\begin{question}{8}
국토의 계획 및 이용에 관한 법령상 개발밀도관리구역에 관한 설명으로 옳지 않은 것은?
\begin{choices}
  \item 개발밀도관리구역의 지정권자는 특별시장ㆍ광역시장ㆍ특별자치시장ㆍ특별자치도지사ㆍ시장 또는 군수이다.
  \item 개발밀도관리구역은 기반시설의 설치가 용이한 지역을 대상으로 건폐율ㆍ용적률을 강화하여 적용하기 위해 지정한다.
  \item 개발밀도관리구역에서는 해당 용도지역에 적용되는 용적률 최대한도의 50퍼센트 범위에서 용적률을 강화하여 적용한다.
  \item 지정권자가 개발밀도관리구역을 지정하려면 해당 지방자치단체에 설치된 지방도시계획위원회의 심의를 거쳐야 한다.
  \item 지정권자는 개발밀도관리구역을 지정ㆍ변경한 경우에는 그 사실을 당해 지방자치단체의 공보에 게재하는 방법으로 고시하여야 한다.
\end{choices}
\end{question}
\answer{2}
